\section{The Heavenly Jerusalem on the Altar: Liturgy, Examination, and Received Prayer}

A Catholic does not merely read about heaven or wait passively for death to unveil it. Every morning across the earth, when the priest ascends the altar to offer the Holy Sacrifice of the Mass, the veil between the Church Militant and the Church Triumphant parts, and the heavenly Jerusalem touches the earth.

Before the Canon begins, the priest extends his hands and summons the faithful upward: \textit{Sursum corda}---``Lift up your hearts.'' They answer: \textit{Habemus ad Dominum}---``We lift them up unto the Lord.''\endnote{\textit{Manual of Prayers for the Use of the Catholic Laity}, Prepared and Published by Order of the Third Plenary Council of Baltimore (New York: Catholic Publication Society, 1889), Ordinary of the Holy Mass, pp.\ 135--137.} At once the Preface lifts earthly thanksgiving into the eternal worship of ``Angels and Archangels, Cherubim also and Seraphim, who cease not daily to cry out with one voice,'' breaking into the \textit{Sanctus} heard by Isaias (Isa 6:3) and St.\ John (Apoc 4:8): ``Holy, Holy, Holy, Lord God of Hosts. Heaven and earth are full of Thy glory. Hosanna in the highest.''

Inside the stillness of the Roman Canon, three prayers anchor the earthly altar directly within the court of heaven.\endnote{1604 Vatican typical \textit{Missale Romanum}, Canon Missae (\textit{Communicantes}, \textit{Supplices te rogamus}, \textit{Nobis quoque peccatoribus}), collated with Eugene Cummiskey's bilingual \textit{Roman Missal for the Use of the Laity} (Philadelphia, 1861), pp.\ xxxiv--xxxviii.} Before the Consecration, the \textit{Communicantes} enters openly into fellowship with the blessed:

\begin{quote}
\textit{Communicantes, et memoriam venerantes, in primis gloriosae semper Virginis Mariae, Genitricis Dei et Domini nostri Jesu Christi: sed et beatorum Apostolorum ac Martyrum tuorum, Petri et Pauli \dots\ et omnium Sanctorum tuorum: quorum meritis precibusque concedas, ut in omnibus protectionis tuae muniamur auxilio.}
\end{quote}

``Communicating with, and honouring in the first place, the memory of the ever glorious Virgin Mary, Mother of our Lord and God Jesus Christ: as also of the blessed Apostles and Martyrs \dots\ and of all thy Saints; by whose merits and prayers grant that we may be always defended by the help of thy protection.''\endnote{Eugene Cummiskey, \textit{The Roman Missal, Translated into the English Language for the Use of the Laity} (Philadelphia, 1861), Canon of the Mass, pp.\ xxxiv, xxxvii--xxxviii. Note that the clause naming St.\ Joseph (\textit{sed et beati Joseph, ejusdem Virginis Sponsi}) was inserted into the \textit{Communicantes} by Pope St.\ John XXIII on 13 November 1962, shortly after the promulgation of the 1962 typical edition.}

After the Consecration, bowing low before the risen Body and Blood of Christ upon the altar, the priest prays the \textit{Supplices te rogamus}:

\begin{quote}
\textit{Supplices te rogamus, omnipotens Deus: jube haec perferri per manus sancti Angeli tui in sublime altare tuum, in conspectu divinae majestatis tuae: ut quotquot, ex hac altaris participatione, sacrosanctum Filii tui Corpus et Sanguinem sumpserimus, omni benedictione caelesti et gratia repleamur.}
\end{quote}

``We most humbly beseech thee, Almighty God, command these things to be carried by the hands of thy holy angels to thy altar on high, in the sight of thy divine Majesty, that as many as shall partake of the most sacred body and blood of thy Son at this altar, may be filled with every heavenly grace and blessing.''\endnote{Eugene Cummiskey (1861), pp.\ xxxvii--xxxviii. While the 1604 typical Latin reads the singular \textit{sancti Angeli tui} (``thy holy angel''), the 1861 Cummiskey English translates the phrase in the plural (``thy holy angels'').} Then, after praying that the faithful departed may receive ``a place of refreshment, light, and peace'' (\textit{locum refrigerii, lucis et pacis}), he strikes his breast at \textit{Nobis quoque peccatoribus}:

\begin{quote}
\textit{Nobis quoque peccatoribus famulis tuis, de multitudine miserationum tuarum sperantibus, partem aliquam et societatem donare digneris, cum tuis sanctis Apostolis et Martyribus \dots\ et omnibus Sanctis tuis: intra quorum nos consortium, non aestimator meriti, sed veniae, quaesumus, largitor admitte.}
\end{quote}

``Also to us sinners, thy servants, confiding in the multitude of thy mercies, vouchsafe to grant some part and fellowship with thy holy apostles and martyrs \dots\ and with all thy saints: into whose company we beseech thee to admit us, not in consideration of our merit, but of thy own gratuitous pardon.'' Naming John the Baptist, Stephen, Matthias, Barnabas, Ignatius, Alexander, Marcellinus, Peter, Felicity, Perpetua, Agatha, Lucy, Agnes, Cecilia, and Anastasia as real brothers and sisters awaiting us, the Church begs entry solely through God's lavish pardon.

Each year on 1 November, the Feast of All Saints gathers all these threads into one radiant solemnity.\endnote{1962 \textit{Missale Romanum}, In Festo Omnium Sanctorum (1 November), Introit, Collect, Lesson (Apoc 7:2--12), Gradual (Ps 33:10--11), Alleluia (Matt 11:28), Gospel (Matt 5:1--12), Offertory (Wis 3:1--3), Secret, Communion (Matt 5:8--10), Postcommunion; collated with the 1922 Tours Mame edition and Rev.\ F.\ X.\ Lasance's \textit{New Roman Missal} (1945).} After \textit{Gaudeamus omnes}, the Lesson proclaims St.\ John's vision of the innumerable white-robed throng before the Lamb (Apoc 7:2--12), and the Gospel answers with the Eight Beatitudes (Matt 5:1--12), placing the vision of heaven directly beside the road that leads there and echoing at Communion, ``Blessed are the clean of heart: for they shall see God.'' In the classical pre-1955 Roman Breviary---before the Sacred Congregation of Rites' general decree \textit{Cum nostra hac aetate} (23 March 1955) under Pope Pius XII suppressed the Octave of All Saints---the Second Day within the Octave (2 November, when the festal Office preceded and stood alongside the Commemoration of All Souls) read Venerable Bede's homily on the heavenly fatherland at Matins:

\begin{quote}
O with what a passion of welcome doth the land above receive them that are coming home from the battle, and go forth to meet them that bear the spoils of the downtrodden foe! \dots\ Come now, therefore, my brethren, and let us enter upon the way; let us turn our faces homeward and heavenward, toward that land wherein our names are written down and we are citizens elect.\endnote{\textit{The Roman Breviary}, Summer (Autumn) Volume, Second Day within the Octave of All Saints (2 November, in the pre-1955 calendar prior to the suppression of the Octave by the Sacred Congregation of Rites' decree \textit{Cum nostra hac aetate}, 23 March 1955), Matins Lessons IV--V (from Venerable Bede), trans.\ John, Marquess of Bute (Edinburgh and London: William Blackwood and Sons, 1879), p.\ 1335.}
\end{quote}

And at Sunday Matins and solemn thanksgivings, the Church sings the \textit{Te Deum laudamus}, uniting angels, Cherubim, Seraphim, Apostles, Prophets, and Martyrs in praise of Christ, who ``having overcome the sting of death, didst open to believers the kingdom of heaven.''

\subsection*{Examination of Conscience}

Before God, a believer may test his own heart against the promise of the kingdom:

\begin{itemize}
\item \textit{Theological hope against despair and presumption.} Do I genuinely long for heaven above earthly comfort, or have I settled into exile as though it were home? Have I despaired of God's mercy over my falls, or presumed on heaven while refusing to amend my life?
\item \textit{Purity of heart and disordered attachments.} ``Blessed are the clean of heart: for they shall see God'' (Matt 5:8). What ``vinegar'' of worldliness, envy, resentment, or impurity fills my heart, crowding out desire for God? Do I guard my body as flesh destined for the resurrection of glory?
\item \textit{Charity toward Christ's poor and the works of mercy.} The King's ``Come, ye blessed'' (Matt 25:34--36) rests on feeding the hungry, sheltering the stranger, and visiting the sick and lonely. Does hope of heaven actually loosen my grip on money, time, and ease for my household and neighbor?
\item \textit{Reverence at Holy Mass.} When I hear \textit{Sursum corda}, do I lift my mind to join the angels and saints around the altar, preparing diligently for Holy Communion as the living pledge of future glory?
\item \textit{Communion with Our Lady and the Saints.} Do I honor Mary as Queen of Heaven, ask her daily for final perseverance, and imitate the virtues of the saints whose fellowship I hope to share?
\end{itemize}

\subsection*{Received Prayers}

The Church sets upon our lips prayers tested by centuries of Catholic worship, lifting the heart toward the vision of God.

\subsubsection*{The Hymn \textit{Te Deum laudamus}}

\textbf{Latin (\textit{Manual of Prayers}, 1889).}

\begin{quote}
\textit{Te Deum laudamus: te Dominum confitemur.\\
Te aeternum Patrem omnis terra veneratur.\\
Tibi omnes angeli, tibi caeli, et universae potestates:\\
Tibi cherubim et seraphim, incessabili voce proclamant:\\
Sanctus, sanctus, sanctus, Dominus Deus Sabaoth:\\
Pleni sunt caeli et terra majestatis gloriae tuae.\\
Te gloriosus Apostolorum chorus.\\
Te Prophetarum laudabilis numerus.\\
Te Martyrum candidatus laudat exercitus.\\
Te per orbem terrarum sancta confitetur Ecclesia.\\
Patrem immensae majestatis.\\
Venerandum tuum verum et unicum Filium.\\
Sanctum quoque Paraclitum Spiritum.\\
Tu Rex gloriae, Christe.\\
Tu Patris sempiternus es Filius.\\
Tu ad liberandum suscepturus hominem, non horruisti Virginis uterum.\\
Tu devicto mortis aculeo, aperuisti credentibus regna caelorum.\\
Tu ad dexteram Dei sedes, in gloria Patris.\\
Judex crederis esse venturus.\\
Te ergo quaesumus, tuis famulis subveni, quos pretioso sanguine redemisti.\\
Aeterna fac cum Sanctis tuis, in gloria numerari.\\
Salvum fac populum tuum, Domine, et benedic haereditati tuae.\\
Et rege eos, et extolle illos, usque in aeternum.\\
Per singulos dies benedicimus te.\\
Et laudamus nomen tuum in saeculum, et in saeculum saeculi.\\
Dignare, Domine, die isto, sine peccato nos custodire.\\
Miserere nostri, Domine, miserere nostri.\\
Fiat misericordia tua, Domine, super nos: quemadmodum speravimus in te.\\
In te, Domine, speravi: non confundar in aeternum.}\endnote{Latin and English texts of the \textit{Te Deum laudamus} from \textit{Manual of Prayers for the Use of the Catholic Laity}, Prepared and Published by Order of the Third Plenary Council of Baltimore (New York: Catholic Publication Society, 1889), pp.\ 160--162, resolving nineteenth-century OCR character errors without altering received wording.}
\end{quote}

\textbf{English (\textit{Manual of Prayers}, 1889).}

\begin{quote}
We praise Thee, O God: we acknowledge Thee to be the Lord.\\
All the earth doth worship Thee: the Father everlasting.\\
To Thee all Angels; to Thee the Heavens and all the Powers therein:\\
To Thee the Cherubim and Seraphim; cry with unceasing voice:\\
Holy, Holy, Holy: Lord God of Hosts.\\
The heavens and the earth are full: of the majesty of Thy glory.\\
Thee the glorious choir: of the Apostles,\\
Thee the admirable company: of the Prophets,\\
Thee the white-robed army of Martyrs: praise.\\
Thee the Holy Church throughout all the world: doth acknowledge.\\
The Father: of infinite Majesty.\\
Thine adorable, true: and only Son.\\
Also the Holy Ghost: the Paraclete.\\
Thou art the King of Glory: O Christ.\\
Thou art the everlasting Son: of the Father.\\
Thou having taken upon Thee to deliver man: didst not abhor the Virgin's womb.\\
Thou having overcome the sting of death: didst open to believers the kingdom of heaven.\\
Thou sittest at the right hand of God: in the glory of the Father.\\
We believe that Thou shalt come: to be our Judge.\\
We beseech Thee, therefore, help Thy servants; whom Thou hast redeemed with Thy precious Blood.\\
Make them to be numbered with Thy Saints; in glory everlasting.\\
O Lord, save Thy people: and bless Thine inheritance.\\
Govern them: and lift them up for ever.\\
Day by day: we bless Thee.\\
And we praise Thy name for ever: and world without end.\\
Vouchsafe, O Lord, this day: to keep us without sin.\\
Have mercy on us, O Lord; have mercy on us.\\
Let Thy mercy, O Lord, be upon us: as we have hoped in Thee.\\
O Lord, in Thee have I hoped: let me never be confounded.
\end{quote}

\subsubsection*{Collect of the Feast of All Saints}

\textbf{Latin (1962 \textit{Missale Romanum}).}

\begin{quote}
\textit{Omnipotens sempiterne Deus, qui nos omnium Sanctorum tuorum merita sub una tribuisti celebritate venerari: quaesumus; ut desideratam nobis tuae propitiationis abundantiam, multiplicatis intercessoribus, largiaris. Per Dominum.}\endnote{1962 \textit{Missale Romanum}, In Festo Omnium Sanctorum (1 November), Collect, collated against the 1922 Tours Mame typical edition.}
\end{quote}

\textbf{English (\textit{Manual of Prayers}, 1889).}

\begin{quote}
Almighty, Everlasting God, who givest us to venerate in one Solemnity the merits of all Thy saints: we beseech Thee that through the multitude of our intercessors Thou wouldst bestow upon us the fulness of Thy mercy, which we most humbly desire. Through our Lord Jesus Christ, etc.\ Amen.\endnote{\textit{Manual of Prayers for the Use of the Catholic Laity} (New York: Catholic Publication Society, 1889), Collect for All Saints, p.\ 773.}
\end{quote}

\subsubsection*{Prayer of St.\ Alphonsus Maria de Liguori}

\begin{quote}
My dear Jesus, you have taught me this prayer,---thy kingdom come. I therefore now pray that your kingdom may come into my heart, so that you may possess it entirely, and that my soul may possess you, her sovereign good. O my Jesus, you have spared nothing in order to save me and to acquire my love: save me then, and let my salvation consist in loving you always in this life and in the next. I have so often turned my back upon you, and yet you assure me that you will not disdain to embrace me for all eternity in Heaven, with as much love as if I had never offended you. And knowing this, can I ever love any thing else but you, who wish to bestow Heaven upon me after I had so often deserved Hell? Ah my Lord, would that I had never offended you! Oh! were I born again, I would always love you. But what is done, is done. I can now do no more than give you the remainder of my life. Yes, I give it entirely to you. I consecrate myself entirely to your love. Depart from my heart, all earthly affections; give place to my God, who wishes to possess it entirely. Yes; take possession of my whole being, my Redeemer, my love, and my God. From this day forward I will think only of pleasing you. Assist me by your grace. I hope in your merits. Augment in me continually your love, and the desire of pleasing you. Heaven! Heaven! When, O Lord, shall the day come, on which I shall see you face to face, and embrace you, without fear of ever more losing you? Ah my God, keep your hand upon me, that I may never more offend you. O Mary, when shall I see myself at your feet in Heaven? Assist me, my mother; do not permit me to be lost, and to go for ever to a distance from you and your Son.\endnote{St.\ Alphonsus Maria de Liguori, \textit{Preparation for Death}, Consideration XXIX (``On Heaven''), Third Point, ``Affections and Prayers,'' trans.\ by a Catholic clergyman (Dublin: James Duffy, 1843), pp.\ 361--362.}
\end{quote}
